% !PRINTED 347
\weberpage{347}
%<<CONT
so ist $P_0$ nicht mehr transitiv, d.\,h. nach Adjunction von $x_0$
zerfällt die Gleichung für die übrigen acht Wurzeln in mehrere Factoren. Wenn die
Gruppe $P$ überhaupt transitiv ist, so kann hierbei $x_0$ jede beliebige der Wurzeln
sein; denn wenn $\pi$ eine Permutation von $P$ ist, durch die $x_0$ in $x_1$ übergeht,
so bleibt durch $P_1 = \pi^{-1} P_0 \pi$ die Wurzel $x_1$ ungeändert, und wenn $P_0$
für acht Wurzeln intransitiv ist, so muss auch $P_1$ für acht Wurzeln intransitiv sein,
da sich die Permutationen von $P_1$ von denen von $P_0$ nur durch die Bezeichnung der
Wurzeln unterscheiden (Bd. I, §.\,154, 6.).

Die einfach transitiven unter den linearen Gruppen haben einen specielleren Charakter
als die Gruppe der Tripelgleichungen.

Wir haben im §.\,78 eine Zerlegung der allgemeinen linearen Congruenzgruppe für den
Modul 3 kennen gelernt, an der wir diese Verhältnisse übersehen können.

Die \gesperrt{Abel'sche} Gruppe $Q$, die aus den Substitutionen

\weq{9}{\xi' \equiv \xi + \alpha, \quad \eta' \equiv \eta + \beta \quad (\mathrm{mod}\ 3)}

besteht, ist gewiss nur einfach transitiv; denn wenn eine Wurzel durch (9) ungeändert
bleiben soll, so muss $\alpha \equiv 0$, $\beta \equiv 0$ sein und alle Wurzeln bleiben
ungeändert. Jede Wurzel ist rational durch jede andere ausdrückbar.

Suchen wir nun in der Gruppe $P$ die Gruppe $P_0$ der Substitutionen auf, durch die
$(0,0)$ ungeändert bleibt, so enthält $P_0$ nur homogene Substitutionen und ist also
der grösste gemeinschaftliche Theiler von $P$ und der homogenen Congruenzgruppe $S$,
und es ist $P = P_0 Q$.

Bilden wir für $S$ die Compositionsreihe (§.\,78):

\[
S,\ S_1,\ S_2,\ S_3,\ S_4,
\]

so ist $S_4 = \begin{pmatrix} 1, & 0 \\ 0, & 1 \end{pmatrix}, \begin{pmatrix} -1, & 0 \\ 0, & -1 \end{pmatrix}$,
und durch $S_4$ gehen die Wurzeln $(1,1)$, $(2,2)$ nur in einander über. Durch $S_4 Q$
kann das Tripel $(0,0)$ $(1,1)$ $(2,2)$ zwar jede Anordnung seiner Wurzeln erfahren.
Es kann aber nur noch in die beiden anderen Tripel $(1,0)$ $(2,1)$ $(0,2)$;
$(0,1)$ $(1,2)$ $(2,0)$ übergehen. $S_4 Q$ ist also nur \gesperrt{einfach transitiv} und
ist überdies \gesperrt{imprimitiv}.

Durch die Gruppe $S_3$, die aus $S_4$ durch Hinzunahme der
Substitution

\[
\begin{pmatrix} -1, & 1 \\ 1, & 1 \end{pmatrix}
\]

%CONT>>
